\section{Method}
\label{sec:method}

WarmCast runs as a controller next to the platform scheduler. Time is divided
into intervals of length $\tau$, which we set to one minute. At the start of
each interval the controller updates a rate estimate for every function,
computes a target pool size, and starts or stops pre-warmed instances to
reach it. Requests are routed to warm instances first; a request that finds
none triggers a cold start as before.

\paragraph{Rate estimation.}
Let $x_t$ be the number of requests a function received in interval $t$. We
estimate its arrival rate with exponential smoothing,
\begin{equation}
  \hat{\lambda}_t = \alpha \, \frac{x_t}{\tau} + (1 - \alpha) \, \hat{\lambda}_{t-1},
  \label{eq:ewma}
\end{equation}
where $\alpha \in (0, 1]$ controls how quickly the estimate follows a change
in load. A large $\alpha$ reacts to bursts but also to noise; we use
$\alpha = 0.3$, which performed best on a held-out day of the trace.

\paragraph{Pool sizing.}
Given the rate, we model the number of arrivals $N$ in the next interval as a
Poisson variable with mean $\mu = \hat{\lambda}_t \tau$. If $k$ instances are
warm, every arrival beyond the $k$-th concurrent one causes a cold start. Let
$c_m$ be the cost of keeping one instance in memory for an interval and $c_s$
the cost of one cold start. The expected cost of a pool of size $k$ is
\begin{equation}
  C(k) = c_m \tau \, k + c_s \, \mathbb{E}\big[(N - k)^{+}\big],
  \label{eq:cost}
\end{equation}
and the controller chooses $k^{*} = \arg\min_{k \ge 0} C(k)$. Because $C$ is
convex in $k$, the minimiser is the smallest $k$ for which
$\Pr[N > k] \le c_m / c_s$, which we find by evaluating the Poisson tail from
$k = 0$ upward.

\paragraph{Cost parameters.}
We express both costs in GB-seconds of memory. The memory cost $c_m$ is the
memory size of the function times $\tau$, and the start cost $c_s$ is the
memory size times the measured cold-start duration, multiplied by a weight
$w$ that encodes how much the operator values latency over memory. With
$w = 1$ the controller minimises total memory time; larger values buy lower
latency with more memory.

\paragraph{Safety limits.}
To avoid holding memory for a function whose traffic has stopped, the
controller caps $k^{*}$ at the largest number of concurrent requests observed
in the last hour and releases all instances of a function that received no
request for 30 minutes. Pre-warming uses spare capacity only; when a host
runs short of memory, pre-warmed instances are evicted before instances that
serve requests.
